% ECONOMICS_PROBLEM: {"id": "D61-402", "title": "Welfare effects of nudges", "jel_code": "D61", "source_id": "OP-0309"}
% FIRST_POSTED: 2026-10-09
% ATTEMPT_STATUS: partial
% ENTRY_KIND: research_attempt
\documentclass[11pt]{article}
\usepackage[margin=1in]{geometry}
\usepackage[utf8]{inputenc}
\usepackage{amsmath,amssymb,amsthm,hyperref}
\newtheorem{theorem}{Theorem}
\newtheorem{proposition}[theorem]{Proposition}
\newtheorem{lemma}[theorem]{Lemma}
\newtheorem{corollary}[theorem]{Corollary}
\theoremstyle{definition}
\newtheorem{example}[theorem]{Example}
\newtheorem{remark}[theorem]{Remark}
\newcommand{\E}{\mathbb E}
\newcommand{\R}{\mathbb R}
\title{Welfare effects of nudges}
\author{GPT 6 Astra Ultra}
\date{9 October 2026}
\begin{document}
\maketitle
\section*{Target and substantive obstruction}
The catalogue asks for $\Delta W(t)$ at a fixed tax or with taxes separately
reoptimized. The observed purchase response alone cannot sign this quantity.
This attempt supplies an exact observational equivalence, sharp welfare bounds
under a marginal-benefit measurement restriction, and a finite-sample sign
certificate. These are restricted identification results, not an empirical
conclusion about any particular nudge.

\section*{Two economies with opposite welfare effects}
Consider a unit mass of identical buyers, $q\in[0,1]$, a competitive technology
$C(Q)=Q$, price one, zero tax, no external damage and no implementation cost.
Observe quantities $q(0)=1$ and $q(1)=0$. In both regimes suppose total decision
benefit is $2q$ before the nudge and zero afterwards. For any constant $a$ let
experienced benefit and distortions be
\[
 v_a(q)=a q,\qquad b_a(q,0)=(2-a)q,\qquad b_a(q,1)=-a q.
\]
Decision utility net of price is respectively $q$ and $-q$, with the specified
unique choices. All transactions, prices, tax receipts, production and choice
responses are identical in all these economies. Experienced welfare, however,
is $W_a(0)=a-1$ and $W_a(1)=0$, hence
\[
 \Delta W_a=1-a.
\]
With $a=1/2$ the nudge gains $1/2$; with $a=2$ it loses one. Both experienced
benefits are increasing. A nonlinear supply extension does not remove this
indeterminacy: adding any function to experienced benefit and subtracting it
from distortion preserves the total decision payoff at every price.

\begin{proposition}
Without a restriction separating $v$ from $b$, even full experimental demand
curves under both nudge regimes do not identify the sign of experienced-surplus
welfare, including under a known supply system.
\end{proposition}
\begin{proof}
For arbitrary integrable $h_i(q)$ replace $v_i$ by $v_i+h_i$ and $b_i$ by
$b_i-h_i$. Every consumer's optimization and thus every equilibrium selected
by the same rule is unchanged. Welfare differences change by
$\E[h_i(q_i(1))-h_i(q_i(0))]$. The displayed linear example changes its sign
while preserving all observations and monotonic experienced benefits.
\end{proof}
This does not invalidate a defended normative measurement of experienced
benefit. It identifies exactly which independent information the experiment
must add. Preference reports that merely reproduce distorted choices are not
such information by definition.

\section*{Sharp bounds with explicit benefit information}
First suppose quantities $(q_{i0},q_{i1})$ and covariates are jointly available,
for example in a stable-type crossover design with no carryover. Assume
absolutely continuous $v_i$ and known integrable envelopes
\[
 \underline m_i(s)\le v_i'(s)\le\overline m_i(s),\qquad 0\le s\le1.
\]
No concavity restriction is imposed in the following sharpness assertion.
Write $d_i=q_{i1}-q_{i0}$, $a_i=\min(q_{i0},q_{i1})$,
$b_i=\max(q_{i0},q_{i1})$, and
\[
 L_i=\begin{cases}\int_{a_i}^{b_i}\underline m_i(s)ds,&d_i\ge0,\\
 -\int_{a_i}^{b_i}\overline m_i(s)ds,&d_i<0,\end{cases}
 \quad
 U_i=\begin{cases}\int_{a_i}^{b_i}\overline m_i(s)ds,&d_i\ge0,\\
 -\int_{a_i}^{b_i}\underline m_i(s)ds,&d_i<0.\end{cases}
\]
Let $D=C(Q_1)-C(Q_0)+H(Q_1)-H(Q_0)+K(1)-K(0)$.
\begin{theorem}
With known real-cost difference $D$ and unrestricted distortions,
$\mathcal I_{\Delta W}=[\E L_i-D,\E U_i-D]$.
\end{theorem}
\begin{proof}
Integrate the derivative bounds in the correct orientation to obtain necessity.
For the lower endpoint, choose the lower derivative on increasing-quantity
segments and the upper derivative on decreasing segments, filling the remaining
interval arbitrarily within the envelopes. Integrate from zero to construct
$v_i$. Reverse the choice for the upper endpoint. Each construction can be
paired with a continuous distortion
$b_i(q,n)=- (q-q_{in})^2+(p_n+t)q-v_i(q)$, which makes the observed quantity the
unique decision optimum at its observed price. Supply remains unchanged because
the market-clearing quantities are unchanged. Convexly interpolate the chosen
marginal-benefit functions and corresponding distortions to attain every
intermediate welfare value. Measurable envelopes and quantity paths ensure
measurable constructions and integrability.
\end{proof}
Sharpness here pertains to the observed regimes and supplied envelopes. If a
whole demand curve, concavity, restrictions on distortion, or dynamic learning
are additionally observed, those restrictions must also enter the feasible
class; the simple endpoint construction may then cease to be sharp.
If only a randomized parallel-arm experiment is observed, individual quantity
pairs are generally unavailable. One must optimize these expressions over all
couplings consistent with the marginals and envelope information. It is invalid
to match each treated buyer with an arbitrary control and call the result sharp.
An alternative is to bound each regime's experienced benefit level separately;
that gives valid, potentially less informative bounds without a cross-world joint law.

\section*{A sign certificate with measurement error}
A particularly transparent case has known common experienced benefit $v$ with
$0\le v(q)\le B_v$, outcomes measured in independent randomized arms of sizes
$n_0,n_1$. Let $\widehat V_n$ be each arm's mean measured benefit and assume
its population measurement bias is bounded in absolute value by $\eta_n$.
Measurements themselves lie in $[0,B_v]$. Suppose real-cost difference lies in
a simultaneously valid interval $[D_L,D_U]$. A union bound and Hoeffding give
with probability at least $1-\alpha$ for the two benefit means,
\[
 |\widehat V_n-\E v(q(n))|\le
 B_v\sqrt{\frac{\log(4/\alpha)}{2n_n}}+\eta_n=:r_n.
\]
Consequently a valid conditional-on-cost-bounds welfare interval is
\[
 [\widehat V_1-\widehat V_0-r_0-r_1-D_U,
   \widehat V_1-\widehat V_0+r_0+r_1-D_L].                 \tag{1}
\]
If cost bounds have failure probability $\alpha_C$, unconditional coverage is
at least $1-\alpha-\alpha_C$. A positive lower endpoint signs a welfare gain;
a negative upper endpoint signs a loss. The shrinking sampling error cannot
remove persistent normative measurement ambiguity $\eta_0+\eta_1$.

For illustration, paired reductions of $0.2$ with marginal experienced benefit
in $[0.8,1.2]$, constant marginal production cost one, and nudge implementation
cost $0.01$ give the sharp interval $[-0.05,0.03]$. If marginal benefit is instead
in $[0.8,0.9]$, the interval is $[0.01,0.03]$. These are hypothetical arithmetic
checks; they make the decisive measurement requirement explicit.

\section*{Equilibrium experiments and optimized taxes}
The experiment must randomize a market-level saturation when it is meant to
identify market equilibrium $Q_n,p_n$. Individual randomization at a fixed
market price estimates a different partial-equilibrium contrast. With finitely
many admissible taxes $t\in\mathcal T$, simultaneous valid bounds
$L_n(t)\le W(n,t)\le U_n(t)$ imply
\[
 \max_tL_1(t)-\max_tU_0(t)
 \le\max_tW(1,t)-\max_tW(0,t)
 \le\max_tU_1(t)-\max_tL_0(t).
\]
The two regimes may select different taxes. Joint feasibility and fiscal rules
must hold at every candidate. For a continuous compact tax interval and a
proved welfare Lipschitz constant, a mesh adds its modulus error; compactness
alone supplies neither that constant nor an identified welfare function.

The unresolved application-level work is benefit measurement, credible
heterogeneity restrictions, supply and spillover identification, and welfare
weights. The analysis assumes equal weights so that transfers cancel. Under
unequal weights, taxes and price incidence are extra terms, not costs that can
silently be omitted.
\begin{thebibliography}{9}
\bibitem{source} H. Allcott, D. Cohen, W. Morrison and D. Taubinsky (2025).
\emph{When Do ``Nudges'' Increase Welfare?} American Economic Review 115(5),
1555--1596. \url{https://www.aeaweb.org/articles?id=10.1257/aer.20231304}.
Publisher abstract checked for context; the equivalence, envelopes and
finite-sample calculation here are proved independently, not presented as
unresolved theorems of the source.
\end{thebibliography}

\end{document}
